\documentclass[12pt]{article}
\usepackage[T1]{fontenc}
\usepackage[utf8]{inputenc}
\usepackage{lmodern}
\usepackage{microtype}
\usepackage[a4paper,margin=1in]{geometry}
\usepackage{setspace}
\usepackage[authoryear,round]{natbib}
\setstretch{1.15}
\setlength{\parindent}{1.2em}
\setlength{\parskip}{0.35em}

\title{The Structural Conditions of Governability in Hybrid Human–AI Systems: A Triadic Theory}
\author{Anonymous manuscript for double-anonymized review}
\date{}

\begin{document}
\maketitle



\begin{abstract}


Hybrid human–AI systems are increasingly embedded in institutional decision-making, yet existing approaches to their governance remain conceptually fragmented, oscillating between technical alignment, regulatory compliance, and boundary-setting frameworks. This article advances a structural claim: governability in hybrid systems depends on the sustained integrity of three irreducible conditions — Field, Coherence, and Limit.

Field designates the relational configuration within which agency and authority become intelligible. Coherence refers to the sustained compatibility of interactions across temporal, functional, and normative dimensions. Limit establishes differentiated boundaries of agency, preserving accountability and the possibility of intervention. Together, these elements constitute a minimal triadic architecture of governability.

The article argues that existing dyadic or modular approaches fail to capture the structural interdependence of these conditions. When any one of the three degrades, systems may continue to function technically or comply formally, yet governability erodes. The triadic model is therefore proposed as a minimal structural condition for institutional governance under conditions of hybrid agency.

Beyond conceptual clarification, the framework is translated into governance requirements, including relational visibility, layered coherence assessment, dynamic boundary calibration, and structural diagnostics of rupture modes. By integrating structural theory with institutional design implications, the article offers both a foundational lens for analyzing hybrid systems and a practical architecture for sustaining adaptive governance.

Governability, on this account, is not guaranteed by performance or control, but emerges from the sustained intelligibility of relations, the compatibility of interactions, and the differentiation of agency.

\end{abstract}

\noindent\textbf{Keywords:} governability; hybrid human–AI systems; AI governance; sociotechnical systems; accountability; institutional design

\section{Introduction}


\subsection*{Governability Under Conditions of Hybrid Agency}


Hybrid human–AI systems are increasingly embedded within institutional decision-making across domains such as public administration, healthcare, finance, and security. These systems operate through dense couplings between human judgment, algorithmic mediation, and organizational structures. Yet as their performance capabilities expand, a structural paradox becomes visible: technical efficiency, alignment optimization, and regulatory compliance do not suffice to secure governability.

Systems may scale, coordinate, and optimize while simultaneously becoming more opaque, less contestable, and progressively incapable of sustaining accountable decision-making. This divergence between performance and governability is not a contingent malfunction but a structural signal: performance can scale independently of the conditions that render power accountable. Governability cannot be inferred from output quality, behavioral alignment, or procedural conformity. It depends upon conditions that precede and structure these phenomena.

The central question, therefore, is not how to improve system outputs alone, but what must structurally obtain for hybrid configurations of human, artificial, and institutional agency to remain governable at all. Under conditions of distributed agency and algorithmic mediation, governability becomes a problem of relational intelligibility, dynamic compatibility, and differentiated authority.

Prevailing approaches to AI governance tend to isolate one dimension of this problem. Sociotechnical and interaction-centered models analyze interfaces and coordination mechanisms. Alignment and optimization research focuses on goal specification, performance convergence, and behavioral control. Regulatory and compliance-based frameworks emphasize boundary-setting and formal accountability mechanisms. Each captures a necessary dimension of governance. None, however, articulates the structural architecture within which these dimensions must cohere for governability to remain possible.

The problem addressed here sits at the intersection of philosophy of technology, sociotechnical systems theory, and political theory. As artificial intelligence becomes institutionally embedded, questions traditionally associated with political order—such as the locatability of authority, the sustainability of institutional decision-making, and the attribution of responsibility—reappear within technologically mediated environments. The governance of hybrid human–AI systems therefore cannot be treated solely as a technical alignment problem or a regulatory design challenge. It raises a more fundamental philosophical question concerning the structural conditions under which political authority remains intelligible and accountable within distributed sociotechnical configurations.

This article advances a structural thesis about the conditions of political order under hybrid agency and proposes a triadic structural model of governability in which \textbf{Field, Coherence, and Limit} constitute irreducible conditions of possibility for governing hybrid human–AI systems.

While the triadic conditions of Field, Coherence, and Limit are here articulated as structural requirements for the governability of hybrid human–AI systems, the architecture they describe is not exclusively technical or institutional. At a more fundamental level, these conditions also correspond to the minimal structure through which human agents orient themselves within complex relational environments. To recognize the field in which one is situated, to maintain coherence in action over time, and to operate within discernible limits are not only requirements of system governance but also basic conditions of intelligible human navigation within sociotechnical worlds.

Field designates the relational configuration within which agency, authority, and decision-making become intelligible. Coherence refers to the sustained compatibility of temporal, functional, and normative dynamics across system layers. Limit establishes differentiated boundaries of agency that render responsibility attributable and intervention meaningful. These conditions are mutually constitutive and non-substitutable. The absence of any one does not merely weaken governance; it renders governability structurally unstable. No configuration lacking one of these conditions can sustain governability without implicitly reinstating it under another conceptual register.

The triadic model is not presented as an additional governance framework, but as a minimal structural articulation of the conditions under which governance frameworks themselves can function. It situates current debates in sociotechnical systems theory, AI alignment research, and accountability scholarship within a more fundamental question: not how to govern hybrid systems effectively, but what makes governing them possible in principle.

This article makes four contributions. First, it introduces a unified triadic architecture that integrates relational, dynamic, and boundary conditions into a single structural model of governability. Second, it demonstrates the irreducibility of these conditions and explains why dyadic or modular approaches permit forms of structural drift that remain undetected by performance metrics. Third, it formalizes modes of structural rupture—field collapse, coherence breakdown, and boundary erosion—as analytic categories for diagnosing fragility in hybrid systems. Fourth, it articulates the institutional implications of structural stewardship, specifying how relational visibility, cross-layer compatibility, and boundary articulation must be preserved if governability is to endure.

The remainder of the article proceeds as follows. Section 2 clarifies the structural distinction between performance and governability and situates the argument within existing governance literatures. Sections 3 through 6 develop the triadic conditions—Field, Coherence, and Limit—as irreducible structural components. Section 7 integrates these elements into a unified architecture and analyzes conditions of stability and rupture. Section 8 concludes by articulating the institutional and ontological implications of treating governability not as an outcome of control, but as a structural condition of political order under hybridized agency.

\section{Performance, Control, and the Limits of Fragmented Governance}


Debates on AI governance have developed along distinct intellectual trajectories. Sociotechnical traditions emphasize the co-production of technology and social order, the relational constitution of artifacts, and the situated character of human–machine interaction \citep{Jasanoff2004,Latour2005}. Alignment and control research frames governance primarily as the problem of ensuring that artificial systems pursue specified goals \citep{Bostrom2014,Russell2019}. Normative and accountability scholarship foregrounds responsibility, contestability, and institutional design \citep{Floridi2018,Dignum2019}.

Each tradition captures a necessary dimension of governance. Yet they remain structurally partial.

Interaction-centered approaches illuminate relational processes but rarely theorize the conditions under which those relations remain intelligible and governable across time. Alignment frameworks address behavioral convergence while abstracting from the relational architecture within which authority must remain locatable and accountable. Regulatory models articulate boundaries of responsibility but often presuppose stable relational fields and coherent dynamics that hybrid systems themselves destabilize.

The limitation is therefore not conceptual weakness but structural fragmentation. Interaction, optimization, and boundary-setting are treated as separable problems. What remains under-theorized is the minimal structural grammar within which relational intelligibility, temporal compatibility, and differentiated authority must coexist for governability to persist.

The triadic model advanced here addresses this gap not by adding another framework, but by articulating the structural conditions presupposed by all of them.

\subsection{Canonical Terminology of the Triadic Model}


The framework developed in this article will be referred to as the \textbf{Triadic Theory of Governability (TTG)}.

Within TTG:

\begin{itemize}
\item \textbf{Field} designates the structured relational horizon within which agency and authority become locatable.
\item \textbf{Coherence} designates the temporal compatibility of systemic dynamics across heterogeneous layers.
\item \textbf{Limit} designates the differentiated boundary conditions that render authority attributable and intervention possible.
\end{itemize}

These are not policy instruments or normative aspirations. They are irreducible structural conditions of possibility for governability under distributed and hybridized agency.

The remainder of the article employs this terminology in its strict structural sense.

\section*{Problem Statement}


The rapid integration of artificial intelligence into institutional and societal systems has revealed a widening divergence between performance and governability.

Hybrid systems may operate with high levels of technical efficiency—generating accurate outputs, scaling decision-making, and optimizing processes—while simultaneously becoming more opaque, less contestable, and progressively incapable of sustaining accountable authority.

The problem is not malfunction. It is structural drift.

Performance measures output quality. Governability concerns the conditions under which power remains intelligible, attributable, and revisable. A system may coordinate effectively and comply procedurally while losing the relational clarity and differentiated authority that render governance possible.

Existing governance approaches tend to privilege one structural dimension at the expense of others.

Interaction-centered models examine localized exchanges but frequently presuppose a stable relational environment. Optimization-driven approaches prioritize convergence yet do not ensure compatibility across temporal sequences or differentiation of responsibility. Regulatory frameworks articulate limits, but boundaries detached from relational intelligibility and temporal compatibility risk becoming symbolic or destabilizing.

Across these perspectives, fragmentation obscures the conditions under which governability deteriorates. Hybrid systems may satisfy performance benchmarks and regulatory requirements while progressively eroding the structural grammar that makes authority accountable.

The problem addressed here is therefore structural rather than instrumental. It concerns the minimal structural grammar required for governability to exist at all under conditions of distributed and hybridized agency.

Governability presupposes three irreducible conditions:

\begin{itemize}
\item a structured relational Field within which agency is locatable,
\item a temporally compatible Coherence across systemic dynamics,
\item and differentiated Limits that render authority attributable.
\end{itemize}

Remove any one, and what remains may be operation, control, or compliance—but not governability.

The triadic model developed in the following sections does not propose an additional governance framework. It articulates the minimal structural grammar within which governance frameworks themselves must operate if power is to remain political rather than merely effective.

\section{Field}


Among the three elements of the triadic structure, \textbf{Field} names the first and most fundamental: the relational condition under which agency becomes politically intelligible.

\subsection{Definition of Field}


In hybrid human–AI–institutional systems, \textbf{Field} designates the structured relational configuration within which agency, authority, and decision-making become locatable.

Field is not a physical environment, a technical substrate, nor a mere aggregation of actors. It names the relational order that conditions how action can appear as attributable, interpretable, and consequential.

More precisely, Field refers to the structured horizon within which action can be recognized as belonging to agents, authority can be situated, decisions can register as consequential, and responsibility can be ascribed without indeterminacy.

Field precedes discrete interaction in a structural sense. It constitutes the relational space within which interaction becomes politically meaningful.

Absent Field, events may occur and outputs may be produced. Yet such operations would lack a shared locus of intelligibility. What remains is functional coordination without governable order.

Field is therefore the condition of possibility for politically locatable agency.

\subsection{Field as Structural and Ontological Condition}


Field is treated here not as metaphor but as structural necessity \citep{Winner1986}.

It is structural insofar as it organizes the horizon of possible relations prior to particular acts. It is ontological insofar as agency, authority, and accountability cannot be conceived outside a relational configuration that renders them intelligible.

No human, institutional, or artificial agent operates outside a Field. All agency is exercised within a relational configuration that conditions visibility, influence, and interpretive orientation.

Field is indirectly observable through stabilized relational patterns—durable authority pathways, recurring informational asymmetries, persistent distributions of influence. Yet these manifestations instantiate Field without exhausting it. They are empirical crystallizations of a deeper relational condition.

To intervene without attention to Field is to act upon outcomes while neglecting the relational architecture that renders those outcomes intelligible. Structural drift in governance rarely originates in isolated decisions; it emerges through transformation in the Field itself—its opacity, fragmentation, or reconfiguration.

\subsection{Structural Axes of Field Stabilization}


Field cannot be decomposed into empirical components. Yet its stabilization becomes intelligible along recurrent structural axes.

Field attains determinate structure along recurrent relational axes—topological configuration, normative constraint, informational ordering, temporal structuration, and asymmetrical power distribution. These axes do not compose Field as parts compose a whole. Rather, they mark the dimensions along which relational space becomes structured, stabilized, or destabilized.

Relational topology concerns the configuration of connections through which authority and coordination circulate.

Normative constraint concerns the formal and informal expectations that delimit permissible action.

Informational ordering concerns the distribution of visibility and interpretive framing.

Temporal structuration concerns the rhythms through which decisions propagate and consequences accumulate.

Power distribution concerns the asymmetrical allocation of decision capacity and intervention authority.

At any given moment, the configuration of these axes determines the degree to which agency within a Field remains locatable and politically intelligible.

\subsection{Field Under Hybrid Agency}


Hybridization does not merely increase complexity; it reconstitutes the Field itself.

The institutional embedding of automated processes alters authority pathways. Optimization cycles reshape informational salience. Algorithmic mediation reorganizes temporal rhythms. Human adaptation to these mediations further transforms relational topology.

The introduction of artificial agents does not simply add participants. It reorganizes the relational configuration within which agency is situated.

Field under hybrid conditions is dynamic but not indeterminate. Its evolution is structurally consequential. Governability depends on whether this evolving relational order preserves sufficient intelligibility for agency to remain locatable and responsibility attributable.

\subsection{Field and the Possibility of Governability}


Governability does not originate in isolated acts of control. It presupposes a relational configuration in which agency can be situated without ambiguity.

Where Field retains structural stability, agents can be identified, authority can be located, decisions can be traced, and responsibility can be attributed within a shared horizon.

When Field destabilizes, governability erodes even in the presence of technical performance. Authority pathways become indeterminate. Informational ordering exceeds interpretive recovery. Temporal propagation outpaces institutional comprehension. In such conditions, governance becomes reactive, intervening at the level of behavior while the relational conditions of behavior remain structurally unexamined.

Field therefore marks the spatial condition of governable order.

\subsection{Structural Distortions of Field}


Field may persist while becoming structurally distorted.

Over-constrained Fields collapse plurality into enforced conformity, compressing relational space to the point of rigidity.

Under-differentiated Fields diffuse authority across layers, rendering responsibility indecidable.

Fragmented Fields lack integrative relational order, stabilizing incompatible local logics.

Opaque Fields obscure authority and causal pathways beyond interpretive recovery.

These distortions do not necessarily interrupt technical operation. They erode the relational conditions under which agency can appear as accountable. When stabilized, they prefigure systemic fragility.

\subsection{Field as Necessary but Not Sufficient}


Field is irreducible because it establishes the relational locus within which other structural dimensions can operate.

It is not reducible to Coherence, which concerns the compatibility of dynamics within a Field. Nor is it reducible to Limit, which concerns the differentiation of authority within that relational space.

Field provides the spatial horizon within which agency becomes locatable. Coherence provides temporal compatibility. Limit provides differentiated boundaries.

Field alone does not secure governability. A relational configuration may exist while remaining dynamically incoherent or normatively undifferentiated.

Field is therefore necessary but not sufficient.

\subsection{Transition}


If Field constitutes the spatial condition under which agency becomes locatable,

\textbf{Coherence} concerns the temporal condition under which that relational configuration remains compatible.

The following section formalizes Coherence as a distinct structural dimension that emerges within Field yet cannot be reduced to it.

\section{Coherence}


\subsection*{Coherence as the Structural Condition of Temporal Compatibility}


\subsection{Coherence and the Temporal Condition of Order}


If Field establishes the structured relational horizon within which agency becomes politically locatable, Coherence designates the condition under which that horizon remains temporally compatible.

Political order does not exist merely in space. It unfolds. Decisions propagate. Interpretations sediment. Authority extends across sequences of action. Coherence names the structural condition under which this unfolding does not dissolve into contradiction.

It is not agreement, convergence, or psychological alignment. It is the compatibility of systemic dynamics across time such that actions, interpretations, and institutional processes do not negate one another as they accumulate.

Without Coherence, a relational Field may remain formally structured, yet its internal dynamics fragment. Authority persists in discrete instances, but continuity collapses. Governance becomes episodic rather than sustained.

Coherence is therefore the temporal condition under which governable order can endure.

\subsection{Irreducibility of Coherence}


Coherence cannot be reduced to Field.

A relational configuration may be well-defined while its temporal propagation generates cumulative contradiction. Spatial intelligibility does not guarantee temporal compatibility.

Coherence cannot be reduced to Limit.

Differentiated authority does not ensure that decisions, once enacted, remain mutually interpretable across sequences of interaction. Boundaries may be explicit while systemic dynamics undermine themselves.

Coherence concerns neither spatial location nor boundary articulation, but the compatibility of unfolding processes.

It is the structural condition that prevents systemic self-negation across time.

Where Coherence collapses, what erodes is not immediate operation but continuity of intelligibility. Decisions may still be executed. Outputs may still be produced. Yet the interpretive horizon within which responsibility can be traced becomes unstable.

Coherence is thus categorially distinct within the triadic architecture.

\subsection{Temporal Compatibility as Structural Requirement}


Political order requires more than synchronized action. It requires that successive decisions remain intelligible in relation to one another.

Temporal compatibility entails that interpretive frameworks do not shift so radically as to invalidate prior commitments; that procedural logics do not propagate effects that undermine their own normative premises; and that institutional processes do not accumulate contradictions beyond interpretive repair.

These are not empirical features but structural requirements. A system may function efficiently while silently eroding the compatibility of its own sequences.

Coherence prevents systemic propagation from degenerating into temporal self-negation.

Where this condition fails, fragmentation precedes visible breakdown. The system continues to act, yet no longer integrates its own history.

\subsection{Coherence Under Hybrid Conditions}


Hybridization intensifies the temporal problem of political order.

Human deliberation, institutional procedure, and algorithmic iteration operate under heterogeneous temporalities. Optimization cycles may compress decision intervals beyond institutional reflexivity. Normative commitments may evolve independently of technical updating. Interpretive categories may drift across layers without reciprocal adjustment.

Under such conditions, incoherence does not arise from malfunction but from temporal misalignment among structurally distinct regimes of action.

The danger is not immediate collapse but cumulative contradiction. A system may maintain surface performance while progressively undermining the continuity required for governability.

Coherence functions here as the structural integrator of heterogeneous temporalities.

\subsection{Coherence and the Conditions of Trust}


Trust is not an independent structural dimension. It emerges where Coherence persists.

When systemic dynamics remain compatible across time, expectations stabilize. Interpretive disagreement remains possible without dissolving shared meaning. Errors can be situated as deviations rather than as structural ruptures.

When incoherence stabilizes, trust erodes not because actors withdraw goodwill, but because the temporal continuity necessary for expectation formation collapses.

Trust is therefore derivative of temporal compatibility.

\subsection{Pathologies of Temporal Incompatibility}


Persistent incoherence manifests not as isolated failure but as structural drift in temporal integration.

Interpretive categories may shift without institutional recalibration. Feedback mechanisms may amplify distortion rather than correct it. Decision cycles may desynchronize from oversight capacities. Causal pathways may proliferate without stable attribution.

These are not technical anomalies. They are distortions in the temporal compatibility of relational order.

When such distortions stabilize, governance remains formally operative yet structurally fragile.

\subsection{Necessity Without Sufficiency}


Coherence is necessary because without temporal compatibility, relational intelligibility disintegrates across sequences of action.

Field may remain structured, yet become internally contradictory.

Limits may remain articulated, yet lose stability within shifting interpretive regimes.

Intervention under conditions of incoherence becomes reactive, as no continuous horizon of evaluation persists.

Yet Coherence alone does not secure governability.

Temporal compatibility without relational localization produces abstraction.

Compatibility without differentiated authority produces integration without accountability.

Coherence provides the temporal continuity within which agency remains interpretable.

\subsection{Structural Position Within the Triad}


The triadic structure becomes clearer at the level of temporal analysis.

Field establishes the relational locus within which agency appears.

Coherence sustains the compatibility of that locus across time.

Limit differentiates authority within and across that unfolding.

Field without Coherence yields spatial order without temporal integration.

Coherence without Limit yields compatibility without differentiation.

Limit without Coherence yields rigid segmentation without continuity.

Governability requires the sustained interaction of spatial intelligibility, temporal compatibility, and differentiated authority.

\subsection{Transition}


If Field secures the spatial condition of agency, and Coherence secures its temporal compatibility, then Limit determines whether authority within that unfolding remains differentiated, accountable, and structurally bounded.

The next section formalizes Limit as the third irreducible condition of governability.

\section{Limit}


\subsection*{Limit as the Structural Condition of Differentiated Authority}


\subsection{The Necessity of Limitation Within Relational Order}


If Field establishes the relational horizon within which agency becomes intelligible, and Coherence sustains the compatibility of systemic dynamics across time, Limit determines the conditions under which such compatibility remains governable rather than absorptive.

Unbounded coherence does not produce stability. It produces indistinction. When dynamic alignment extends without structural differentiation, authority fuses with execution, delegation becomes diffusion, and integration collapses into self-referential closure. What appears as seamless coordination conceals the erosion of accountable distinction.

Limitation is therefore not external constraint. It is the structural condition that prevents relational density and dynamic compatibility from dissolving the very distinctions that render governance possible.

Limit renders coherence finite, situated, and revisable.

\subsection{Structural Definition of Limit}


Limit designates the structural condition through which boundaries of agency, interpretation, and authority are differentiated within a relational system.

It does not operate as prohibition imposed from outside the system. It constitutes the differentiation internal to the system through which agency becomes locatable and responsibility becomes attributable.

Without boundary conditions, action may occur and coordination may persist, but authority cannot be meaningfully distinguished from execution, nor responsibility from effect. Limit therefore secures the structural separability of roles within a shared relational horizon.

Limit is categorially distinct from Field and Coherence.

Field provides relational space.

Coherence sustains dynamic compatibility.

Limit differentiates authority within that space and across those dynamics.

\subsection{Limit and the Differentiability of Agency}


Governability presupposes that agency remains differentiable \citep{Ostrom1990}.

To govern is not merely to coordinate behavior, but to sustain distinctions between deciding, executing, supervising, and answering. Where such distinctions erode, responsibility becomes structurally undecidable.

Limit preserves these distinctions by instituting the structural separability of decision, execution, supervision, and answerability.

Without such differentiation, Field may remain relationally structured and Coherence may maintain compatibility, yet governance collapses into either diffusion or absorption. Authority persists, but its location becomes indeterminate.

Limit thus ensures that agency does not dissolve into the very dynamics it helps generate.

\subsection{Limit Under Conditions of Hybrid Agency}


Hybrid human–AI systems intensify the necessity of structural limitation because they amplify asymmetries of speed, scale, and epistemic capacity.

Under such conditions, differentiation is not spontaneously preserved. It must be structurally maintained.

Hybridization intensifies the problem of boundary maintenance because agency becomes operationally distributed without becoming normatively equivalent. To attribute undifferentiated authority within such distributed configurations collapses the distinction between execution and judgment.

Conversely, human and institutional actors may externalize responsibility into automated processes, treating outputs as self-justifying. Such displacement does not eliminate responsibility; it renders it opaque.

Limit functions here not as technological restriction, but as preservation of differentiated agency within distributed configurations of action. It ensures that automation does not silently reconfigure authority beyond the interpretive horizon of accountability.

\subsection{Limit and the Possibility of Accountability}


Accountability presupposes differentiation.

If it is indeterminate whether a decision originates in human judgment, institutional mandate, or algorithmic mediation, responsibility cannot be meaningfully assigned. Appeals lose direction. Correction loses traction.

Limit establishes the structural conditions under which:

\begin{itemize}
\item authority remains locatable,
\item intervention remains viable,
\item responsibility remains attributable,
\item revision remains possible.
\end{itemize}

In its absence, governance persists only symbolically. Rules may exist, but they no longer delimit power. Compliance may be recorded, yet accountability dissolves into procedural form without structural force.

\subsection{Dynamic Limits}


Limits need not be static, but their revision must itself remain structured.

Under systemic evolution, boundary conditions must remain structurally recalibrable without dissolving their differentiating function. Yet revision without procedural coherence reintroduces indeterminacy at a higher level.

Dynamic limitation therefore requires not fluidity, but structured revisability. Boundaries must remain both differentiating and intelligible across time.

Flexibility without differentiation produces diffusion. Rigidity without relational integration produces formalism. Limit sustains the middle condition in which authority remains bounded yet responsive.

\subsection{Pathologies of Limit Failure}


When Limit erodes, agency diffuses. Authority expands without articulation. Responsibility becomes structurally indecidable.

When Limit hypertrophies, differentiation hardens into rigid segmentation. Boundaries proliferate without integration. Constraint replaces intelligibility.

Both conditions undermine governability.

Erosion collapses differentiation into absorption.

Hypertrophy collapses differentiation into formalism.

Limit fails in both cases because it ceases to function as a generative structural boundary.

\subsection{Irreducibility of Limit}


Limit cannot be reduced to Field.

A relational configuration, however structured, does not by itself differentiate authority.

Limit cannot be reduced to Coherence.

Dynamic compatibility, however sustained, does not establish boundaries of responsibility.

A system may exhibit relational intelligibility and dynamic alignment while losing the differentiability required for accountable governance.

Limit is therefore irreducible within the triadic architecture.

Without it, Field becomes expansively indeterminate and Coherence becomes self-referential integration.

Limit provides the differentiating structure within which authority remains attributable.

\subsection{Completion of the Triad}


With Limit, the triadic architecture attains structural closure.

Field establishes the relational horizon of agency.

Coherence sustains its temporal compatibility.

Limit differentiates authority within and across those dynamics.

The absence of any one produces not imperfection, but structural collapse of governability.

\subsection{Concluding Claim}


Limit is not an obstacle to coherence.

It is the condition that prevents coherence from becoming indistinction.

In hybrid human–AI systems, to articulate Limits is not to constrain progress. It is to preserve the differentiability of agency upon which responsible governance depends.

Where Limit is sustained, authority remains political.

Where it erodes, governance becomes either absorption or simulation.

Field makes agency locatable.
Coherence makes it sustainable.
Limit makes it answerable.

\section{The Triadic Architecture of Governability}


\subsection{From Analytical Distinction to Structural Necessity}


Field, Coherence, and Limit were introduced as analytically distinguishable dimensions of governability. This section advances a stronger claim: Field, Coherence, and Limit are not merely complementary lenses, but mutually constitutive and jointly necessary structural conditions. Under conditions of distributed and hybrid agency, governability itself depends upon the preservation of this triadic structure.

Governability does not arise from the aggregation of relational space, dynamic alignment, and bounded authority as separable features. It emerges only where these dimensions operate together as a unified architecture. Field provides no governable order without Coherence; Coherence cannot stabilize without Limit; Limit has no intelligible force outside a structured Field. Each presupposes the others as conditions of its own operability.

The triadic model therefore exceeds analytical pluralism. It articulates a \textbf{structural configuration} within which the exercise of power remains intelligible, attributable, and revisable over time. The relation among the three elements is not additive but constitutive: removing any one transforms the character of the whole.

This interdependence grounds a further claim. If governability requires a structured relational locus, sustained dynamic compatibility, and differentiated authority, then Field, Coherence, and Limit cannot be understood as contingent features of governance arrangements. They correspond instead to structural conditions of political order under hybrid agency.

The argument that follows does not defend the usefulness of the triad. It establishes its \textbf{irreducibility and minimal necessity}.

\subsection{The Three Structural Dimensions of Political Order}


If governability constitutes a structural condition of political order under hybrid agency, then the triadic architecture corresponds to three irreducible structural dimensions: relational intelligibility, temporal compatibility, and differentiated authority. These are not empirical variables but categorial requirements.

\subsubsection{Field and Relational Intelligibility}


Political order presupposes a structured relational horizon within which action can appear as attributable and authority as locatable \citep{Arendt1958}. Field names this horizon.

Governance is not the mere production of coordinated behavior; it is the exercise of authority within a shared relational space. Without such a space, power may circulate, decisions may be executed, and outputs may be optimized, yet agency lacks a stable locus. Actions occur, but they do not take place within an intelligible configuration of responsibility.

Relational intelligibility is therefore a structural condition. It ensures that agents are not abstract nodes in a network but situated actors within a politically meaningful order. Field provides the configuration within which authority can be recognized as exercised rather than merely enacted.

Where Field degrades, governance collapses into dispersed operation. Authority becomes disembedded from relational orientation. What persists may be coordination, but not governable order.

\subsubsection{Coherence and Temporal Compatibility}


Political order unfolds across time \citep{MarchOlsen1989}. Decisions propagate, interpretations accumulate, and consequences reverberate through institutional and technical layers. Coherence names the condition under which this unfolding remains compatible rather than self-negating.

Temporal sustainability does not require uniformity or consensus. It requires that systemic dynamics do not undermine their own intelligibility as they evolve. Without Coherence, authority may remain formally defined and boundaries may remain intact, yet decisions silently contradict one another. Fragmentation replaces continuity.

Coherence therefore functions as the structural condition of temporal durability. It ensures that relational configurations do not dissolve into episodic alignment or cumulative contradiction. Governance becomes possible not because deviation is eliminated, but because deviation remains interpretable within a stable dynamic horizon.

Where Coherence erodes, governance persists only in fragments. Accountability becomes reactive rather than sustained, and political order loses its continuity.

\subsubsection{Limit and Differentiable Authority}


Political order presupposes distinction. To govern is to exercise authority in a manner that differentiates roles, responsibilities, and capacities for intervention. Limit names the structural condition under which such differentiation remains possible.

Without differentiated authority, agency fuses or diffuses. Human judgment may be absorbed into automation; institutional responsibility may become indeterminate; control may persist without answerability. In such cases, power operates, but it does not remain accountable.

Limit ensures that authority remains locatable and revisable. It establishes the structural boundaries that prevent coherence from collapsing into closure and prevent relational density from dissolving into indistinction. Differentiation is not a restriction on governance; it is its enabling condition.

Where Limit erodes, governance gives way to either uncontrolled expansion or rigid formalism. In both cases, what disappears is the possibility of attributing responsibility within an intelligible order.

\subsection{Dynamic Interdependence}


Field, Coherence, and Limit are analytically distinguishable but structurally inseparable. Their significance emerges only through their dynamic interrelation. The triadic architecture is not hierarchical or sequential; it is internally relational. Each element conditions and is conditioned by the others.

\subsubsection{Field–Coherence}


Coherence presupposes Field because compatibility can only arise within a determinate relational horizon. Without a structured locus of agency, dynamic alignment lacks orientation. There is nothing for compatibility to stabilize.

Conversely, Field without Coherence cannot sustain itself over time. A relational configuration that does not maintain dynamic compatibility fragments into incompatible local logics. What was once an intelligible order becomes a disjointed aggregation of interactions.

Field provides the spatial intelligibility of agency; Coherence sustains its temporal continuity. The absence of either destabilizes both.

\subsubsection{Coherence–Limit}


Coherence requires Limit to avoid self-closure. Dynamic compatibility, if unbounded, risks absorbing differentiation into enforced uniformity. Alignment becomes indistinguishable from control when no structural boundaries preserve plural agency and revisability.

Limit, however, depends on Coherence to remain intelligible. Boundaries that operate within incoherent dynamics become arbitrary or symbolic. Differentiation loses meaning when decisions propagate without compatibility across time and function.

Coherence without Limit yields undifferentiated integration.

Limit without Coherence yields rigid formalism.

Their interdependence preserves both adaptability and accountability.

\subsubsection{Field–Limit}


Limit operates within Field and simultaneously configures it. Boundaries are not external impositions but relational differentiations instituted within a structured space of agency.

Without Field, Limit lacks a locus of application. There is no shared relational configuration within which authority can be differentiated. Boundaries become abstract prescriptions detached from lived relational order.

Without Limit, Field dissolves into indistinction. Relational density expands without differentiation, and authority loses locatability. Agency becomes either fused or diffusely distributed.

Field makes differentiation possible; Limit makes relational order governable.

\subsection{Conditions of Stability}


The dynamic interdependence of Field, Coherence, and Limit does not guarantee durability. Governability persists only when certain structural conditions are maintained. Stability, in this context, does not signify equilibrium or absence of conflict, but the sustained viability of the triadic architecture under conditions of change.

Three principles delineate the minimal requirements for such stability.

\subsubsection{Minimal Coherence Thresholds}


For a relational order to remain governable, its internal dynamics must not fall below a threshold of compatibility. Coherence need not be maximal, but it must be sufficient to prevent systemic self-negation.

Below this threshold, decisions propagate without mutual interpretability. Temporal sequences undermine one another. Functional layers diverge into incompatible logics. Authority may remain formally bounded, yet its exercise becomes episodic and reactive.

Minimal coherence ensures that relational intelligibility is not continually dissolved by the system’s own operations. It preserves the temporal continuity required for accountability to remain structurally sustainable rather than contingent.

\subsubsection{Legible Relational Field}


Field must remain sufficiently intelligible to those operating within it. A relational configuration that becomes opaque to its participants loses its capacity to sustain accountable agency.

Legibility does not require total transparency. It requires that authority pathways, decision trajectories, and asymmetries of influence remain structurally interpretable. Where relational structures become inaccessible or inscrutable, coherence assessments lose grounding and limits lose practical force.

A Field that cannot be situated cannot be governed. Stability therefore presupposes a relational environment that remains intelligible enough to orient action and contestation.

\subsubsection{Operationally Integrated Limits}


Limits must be more than formally articulated; they must be structurally operative within the relational and dynamic conditions of the system.

Boundaries that are detached from the Field lack legitimacy. Boundaries that are misaligned with systemic dynamics lose enforceability. In either case, differentiation becomes symbolic rather than effective.

Operational integration requires that Limits remain compatible with both relational configuration and dynamic propagation. Only then can authority remain locatable, responsibility attributable, and intervention viable without destabilizing the architecture itself.

\subsection{Structural Modes of Rupture}


If the triadic architecture specifies the minimal conditions of governability, its rupture follows from the violation of those conditions. Breakdown is not accidental. It is the structural consequence of failing to sustain relational intelligibility, dynamic compatibility, or differentiated authority.

Each element of the triad admits a corresponding mode of collapse.

\subsubsection{Field Collapse}


When the relational configuration ceases to remain legible, the Field collapses.

This occurs not because interaction stops, but because the structured horizon within which agency is situated disintegrates. Authority pathways become indeterminate. Decision trajectories lose interpretive anchoring. Power continues to operate, yet no longer within a shared relational map.

Under such conditions, actions may still be executed and constraints enforced, but they no longer appear within a politically intelligible space. Coordination persists as technical operation; governability dissolves as relational order.

Field collapse is thus the logical consequence of sustained relational opacity.

\subsubsection{Coherence Breakdown}


When systemic dynamics fall below minimal thresholds of compatibility, Coherence breaks down.

Decisions propagate without temporal alignment. Functional layers generate mutually undermining effects. Normative commitments decouple from operational behavior. The system does not cease functioning; it ceases integrating.

Incoherence does not immediately eliminate structure. Rather, it fragments continuity. Authority remains bounded, yet its exercise produces contradictions that accumulate across time and function.

Coherence breakdown is therefore the structural consequence of incompatible propagation within an otherwise existing Field.

\subsubsection{Limit Erosion and Hypertrophy}


When boundaries cease to be operationally integrated, Limit fails.

Failure may take two opposed but structurally related forms.

\textbf{Erosion} occurs when differentiation weakens. Authority diffuses, responsibility becomes unlocatable, and intervention loses viability. Agency fuses into opacity. Governance persists only nominally.

\textbf{Hypertrophy} occurs when boundaries proliferate or rigidify without relational or dynamic integration. Differentiation hardens into formalism. Constraint substitutes for intelligibility. The system remains bounded, yet loses adaptive viability.

In both cases, Limits cease to function as generative conditions of accountable authority.

Limit failure is thus the structural consequence of boundary detachment — whether through dilution or overextension.


\subsection{Minimality and Categorial Irreducibility}


\subsubsection{Governability as a Condition of Possibility}


Governability is not equivalent to coordination, efficiency, compliance, or control. It designates the structural condition under which collective power can be exercised in a manner that remains intelligible, attributable, and revisable over time.

To govern is not merely to produce outcomes, but to exercise authority within a relational order in which agency can be located, responsibility can be assigned, and intervention remains possible without collapsing the system’s intelligibility.

Under conditions of hybridized agency—where human, institutional, and artificial capacities are densely interwoven—these structural requirements do not diminish. They become more explicit. Governability is therefore not a contingent feature added to political order. It is the structural condition under which political order can exist as such.

The triadic architecture articulates this condition.

\subsubsection{Three Categorial Requirements}


Any system that claims to be governable must satisfy three irreducible requirements.

First, there must exist a structured relational locus within which actions can appear as situated. Without such a Field, power may operate, but it lacks a shared horizon of intelligibility.

Second, systemic dynamics must remain sufficiently compatible across time and function to prevent self-negation. Without Coherence, authority may be bounded, yet its exercise becomes internally contradictory and temporally unsustainable.

Third, agency must remain differentiable. Without Limit, distinctions between deciding, executing, supervising, and answering collapse. Power may persist, but it becomes indistinguishable from absorption or diffusion.

These are not empirical variables. They are structural dimensions of political order. They articulate the minimal grammar within which political power can appear as accountable rather than merely effective.

\subsubsection{Non-Substitutability}


The three dimensions are categorially distinct and therefore non-substitutable.

Relational intelligibility cannot be derived from dynamic compatibility.

Compatibility cannot be inferred from bounded roles.

Bounded roles cannot emerge from relational density alone.

No dyadic configuration can satisfy the structural requirements of governability without covertly presupposing the absent third dimension.

Field and Coherence without Limit yield absorption.

Coherence and Limit without Field yield formalism without relational anchoring.

Field and Limit without Coherence yield fragmentation.

To remove any one element is not to weaken governability incrementally. It is to undermine a categorical condition of its existence.

What remains may be operation, control, or compliance. It is not governability.

\subsubsection{The Impossibility of a Fourth Structural Element}


If the triad is minimal, it must also be exhaustive at the categorial level.

Concepts frequently invoked in governance discourse—transparency, trust, accountability, participation, oversight, power—do not introduce additional structural dimensions. They are expressions, derivatives, or configurations within Field, Coherence, and Limit.

Transparency concerns the legibility of Field and the traceability sustained by Coherence.

Trust emerges from sustained Coherence within stable Limits.

Accountability presupposes differentiated authority (Limit) within an intelligible Field.

Power operates as a relational property of Field, structured and bounded by Limit.

To state the categorial claim in canonical form:

Any purported fourth dimension is analytically reducible to Field, Coherence, or Limit; no independent structural category of governability exists beyond the triad.

The triadic architecture is therefore minimal not by parsimony, but by structural necessity.

To deny one element is not to refine the model.

The triadic architecture does not describe one possible model of governance. It articulates the minimal structural grammar of political order under conditions of distributed and hybridized agency and identifies not how governance ought to be arranged, but what must be structurally present for governance to remain possible at all.

\subsection{Transition to Institutional Translation}


If governability is structurally triadic, institutional design cannot be confined to rule-making, compliance verification, or performance management. It becomes the stewardship of the conditions that render governing possible in the first place.

Institutions do not add a fourth dimension to the architecture. They operate within it. Their task is not to perfect outcomes, but to preserve relational intelligibility, sustain dynamic compatibility, and maintain differentiated authority under conditions of systemic evolution.

Institutional translation, therefore, does not move beyond the triadic structure. It assumes responsibility for its maintenance.

The next sections turn from categorial necessity to structural dynamics and institutional stewardship.

\section{Structural Stability and Cascading Rupture}


The triadic architecture specifies not only the conditions of governability, but the structural logic of its deterioration.

Rupture does not originate in isolated malfunction. It emerges when one of the triadic conditions falls below structural viability and begins to destabilize the others. Because Field, Coherence, and Limit are mutually constitutive, failure is rarely confined to a single dimension. It propagates.

\subsection{Field Destabilization}


Field rupture occurs when the relational horizon of agency loses legibility.

Authority pathways become indeterminate. Interpretive orientation fragments. Decisions continue to occur, yet no longer within a shared relational configuration.

The immediate effect is not operational collapse, but erosion of intelligibility. Once relational orientation weakens, coherence assessments lose grounding and limits lose applicability. What begins as opacity in the Field progressively undermines temporal compatibility and boundary differentiation.

Field destabilization therefore initiates structural cascade.

\subsection{Coherence Destabilization}


Coherence rupture arises when systemic dynamics cease to remain mutually compatible across temporal and functional layers.

Decisions propagate without interpretive continuity. Institutional commitments diverge from operational trajectories. Alignment becomes episodic rather than sustained.

The Field may remain structured, and Limits may remain formally articulated, yet their interaction becomes temporally incoherent. Boundaries apply inconsistently; relational orientation shifts unpredictably.

Incoherence introduces cumulative contradiction. It does not immediately destroy structure, but it erodes its durability.

\subsection{Limit Destabilization}


Limit rupture occurs when differentiation of authority loses structural integration.

Erosion produces diffusion: responsibility becomes unlocatable, intervention loses traction, accountability dissolves.

Hypertrophy produces rigidification: boundaries proliferate without relational grounding, suppressing adaptability while preserving formal segmentation.

In both cases, differentiated authority ceases to function as a generative boundary condition. Without stable Limits, Field becomes expansively indeterminate and Coherence becomes absorptive or closed.

\subsection{Cascading Interaction}


Because the triadic conditions are internally relational, rupture propagates.

Field opacity destabilizes coherence by removing interpretive grounding.

Incoherence destabilizes limits by rendering boundaries arbitrary or symbolic.

Boundary failure destabilizes Field by dissolving the differentiability upon which relational intelligibility depends.

Structural fragility therefore accumulates gradually. Systems may retain surface performance while progressively losing triadic integrity.

Collapse, when it appears, is often misdiagnosed as sudden failure. It is the endpoint of accumulated structural drift.

\subsection{Structural Insight}


Governability fails not when performance declines, but when the triadic architecture falls below minimal viability.

The critical threshold is not output quality, but structural integration.

Where Field remains legible, Coherence remains compatible, and Limits remain differentiating, rupture remains containable.

Where one dimension deteriorates without structural recalibration, cascade becomes probable.

\subsection{Triadic Cascade}


Contemporary crises of governance frequently appear as technical failures, legitimacy disputes, or regulatory inadequacies. Yet many such crises share a deeper structural pattern: they arise when one dimension of the triadic architecture—Field, Coherence, or Limit—begins to deteriorate while the system continues to operate.

Three recurrent structural patterns of crisis tend to emerge across contemporary technological and institutional systems.

\subsubsection{Systemic Opacity — Field Degradation}


Crises of systemic opacity arise when the relational configuration within which agency and authority operate becomes difficult to interpret.

In such cases, systems continue to function and decisions continue to be produced, yet the relational horizon that allows participants to situate action becomes increasingly obscure. Authority pathways become difficult to trace, and responsibility becomes harder to attribute within a shared interpretive map.

Under these conditions, disputes over legitimacy frequently emerge not because outcomes are necessarily incorrect, but because the relational configuration within which those outcomes arise is no longer intelligible to those subject to it.

The deterioration concerns the Field: the relational condition that renders agency politically locatable.

\subsubsection{Temporal Contradiction — Coherence Breakdown}


A second structural pattern appears when systemic dynamics cease to remain compatible across time.

Institutions and technological systems increasingly operate across heterogeneous temporal regimes: human deliberation, institutional procedure, and algorithmic iteration. When these temporalities fall out of alignment, decisions may remain individually valid while collectively generating contradiction.

Policies change faster than interpretive frameworks adapt; automated systems evolve faster than oversight structures recalibrate; institutional commitments accumulate without temporal integration.

Such systems may exhibit high operational performance while gradually losing the capacity to integrate their own history. Governance persists episodically rather than as a sustained temporal order.

This pattern reflects deterioration in Coherence: the structural condition of temporal compatibility.

\subsubsection{Diffuse Authority — Limit Erosion}


A third structural pattern emerges when the differentiation of authority weakens within distributed systems.

Hybrid sociotechnical configurations frequently distribute action across multiple agents: human actors, institutional processes, and automated systems. When the boundaries among these forms of agency lose clarity, responsibility becomes structurally difficult to assign.

Decisions may be produced through chains of delegation in which no single locus of authority retains full interpretive or corrective capacity. Responsibility disperses across layers without remaining attributable to a determinate agent.

Governance does not disappear under such conditions. Instead, it becomes symbolically present while structurally diluted.

This pattern corresponds to deterioration in Limit: the structural differentiation that preserves the locatability of authority and the possibility of intervention.

\subsection{Cascade Synthesis}


These three patterns rarely remain isolated.

Opacity in the relational Field undermines the interpretive grounding required to assess coherence. Temporal incoherence destabilizes the practical application of limits. Boundary erosion, in turn, dissolves the relational distinctions upon which Field intelligibility depends.

Contemporary governance crises therefore often appear sudden, yet they typically represent the endpoint of cumulative structural drift across the triadic architecture.

The analytical significance of the triadic model lies precisely in its ability to render these structural patterns visible before collapse becomes evident.

\section{Institutional Translation}


The triadic architecture does not function merely as a theoretical abstraction; it provides a diagnostic orientation for evaluating institutional fragilities in contemporary hybrid systems.


If the previous section identified the structural conditions under which governability deteriorates, the question that follows is institutional: how can these conditions be preserved?

If the triadic structure articulates the minimal architecture of governability, institutional translation does not introduce an additional element. Rather, it concerns how this minimal architecture can be sustained under conditions of systemic evolution.

Governance, under the triadic model, ceases to be the regulation of outputs and becomes the structured preservation of the conditions under which governing remains possible.

Once Field, Coherence, and Limit are recognized as irreducible conditions of governability, governance can no longer be confined to compliance, optimization, or episodic correction. It assumes responsibility for sustaining the relational architecture within which accountability, intelligibility, and revision remain structurally viable.

The decisive question shifts accordingly. It is no longer whether systems perform or conform. It is whether the structural conditions that render governing intelligible endure across transformation.

\subsection{Governance as Structural Stewardship}


Field, Coherence, and Limit do not maintain themselves. Hybrid systems evolve through scale, automation, institutional embedding, and recursive optimization. Structural conditions, once established, are continuously exposed to drift.

Governance therefore becomes a practice of structural stewardship: the sustained preservation of relational intelligibility, temporal compatibility, and differentiated authority under conditions of systemic change.

This reorientation alters the object of governance. Rather than intervening primarily at the level of discrete decisions, governance concerns the architecture within which decisions become attributable, contestable, and correctable. It is not transactional but relational; not reactive but architectural.

The triad does not correspond to instruments but to dimensions of preservation:

Field concerns relational intelligibility.

Coherence concerns temporal compatibility.

Limit concerns differentiated authority.

Institutional forms are therefore to be assessed not by their regulatory density but by their capacity to sustain these structural dimensions without collapse, absorption, or rigidity.

\subsection{Preserving the Relational Field}


To govern the Field is to preserve the conditions under which agency remains politically locatable.

Relational intelligibility exceeds disclosure. Data may circulate abundantly while authority remains opaque. What must endure is the structured possibility of situating action: who acted, under what mandate, through which mediating structures, and with what avenue of intervention.

Institutional forms that sustain Field integrity tend to render authority pathways explicit, maintain traceability between automated processes and institutional mandates, and treat interpretive friction not as noise but as structural signal.

Participation and contestation assume a structural role in this context. They function less as democratic add-ons than as reflexive mechanisms through which the Field becomes capable of self-observation. Where avenues exist for articulating misalignment, relational opacity can be detected before it stabilizes into illegibility.

Field preservation therefore consists in maintaining a legible relational environment under conditions of density, mediation, and distributed agency. Without such legibility, coherence cannot be assessed and limits cannot remain meaningful.

\subsection{Sustaining Coherence Without Collapsing into Control}


If Field secures the locus of agency, governance must also sustain its temporal continuity.

Coherence concerns whether heterogeneous layers of action remain mutually interpretable as they evolve. The danger lies in mistaking coherence for control. Control minimizes deviation. Coherence preserves compatibility across difference.

A system may converge flawlessly toward predefined objectives while eroding the interpretive bridges that render its operations intelligible. Such systems exhibit performance without governability.

Structural coherence is not engineered uniformity but maintained compatibility. It requires that temporal cycles remain survivable across layers, that optimization logics do not silently displace institutional mandates, and that normative commitments do not detach from operational practice.

Contestation here assumes stabilizing significance. Systems incapable of absorbing dissent through interpretive recalibration may achieve short-term consistency while incubating long-term fragility.

Governability persists not because divergence disappears, but because divergence remains intelligible and corrigible within a shared structural horizon.

\subsection{Preserving Limits as Differentiated Authority}


Where Field secures intelligibility and Coherence sustains continuity, Limits preserve the differentiability of authority.

Hybrid systems distribute action across human judgment, institutional allocation, and algorithmic execution. These forms of agency are structurally distinct. When differentiation erodes, responsibility diffuses and accountability dissolves.

Preserving Limits does not primarily mean imposing constraint. It means sustaining explicit articulations of role, scope, and escalation under conditions of technological mediation. Authority must remain locatable; intervention must remain viable; responsibility must remain attributable.

Limits are boundary processes rather than static prohibitions. As systems evolve, authority configurations require recalibration. Yet recalibration itself must remain procedurally bounded, lest silent expansion or rigidification undermine differentiation.

Without Limits, governance collapses either into automated absorption or formalistic rigidity. With them, hybrid systems may evolve while preserving the structural preconditions of answerability.

\subsection{Structural Diagnostics and Drift}


The triadic structure functions not only as architecture but as diagnostic orientation.

Governance traditionally treats failure as episodic breach \citep{Floridi2018,Dignum2019}. Structural governance recognizes fragility as cumulative drift. The relevant question is not merely whether harm has occurred, but whether the conditions that render harm intelligible and correctable remain intact.

Field degradation appears as relational opacity.

Coherence breakdown manifests as temporal desynchronization or cumulative contradiction.

Limit erosion reveals itself in responsibility diffusion or boundary rigidity.

Because the triadic elements are mutually constitutive, deterioration rarely remains isolated. Opacity destabilizes compatibility; incompatibility undermines boundaries; boundary erosion renders the Field unintelligible.

Structural monitoring therefore concerns the continued viability of governability itself. Informal adaptations, friction, and interpretive disputes become early indicators of architectural stress rather than anomalies to be suppressed.

Governance, in this light, becomes anticipatory rather than merely corrective.

\subsection{Toward Structural Political Theory}


The triadic model does not compete among regulatory doctrines. It articulates the structural horizon within which any doctrine must operate if governability is to remain possible.

Governance in hybrid human–AI systems cannot be reduced to output regulation or performance maximization. It becomes the disciplined preservation of relational conditions.

Adaptive governance does not signify regulatory expansion but structural vigilance. Institutions recalibrate not to maximize control, but to sustain the possibility of governing.

Governability is not a fixed achievement. It is a sustained structural accomplishment.

When Field remains intelligible, Coherence remains compatible, and Limits remain differentiating, hybrid systems retain the capacity for correction, contestation, and responsible adaptation. When any element erodes, governance persists only symbolically.

The triadic model therefore intervenes not merely in AI policy debates but in the architecture of political theory itself. It reframes governability as a structural condition composed of relational intelligibility, temporal compatibility, and differentiated authority. In doing so, it offers a minimal ontology of political order under conditions of hybridized agency.

If contemporary crises of legitimacy, opacity, and institutional fragility signal anything, it is not simply regulatory insufficiency but structural drift. The question is no longer how to optimize governance, but how to preserve the conditions under which governance remains possible.

\section{Conclusion}


\subsection*{Governability as a Structural Achievement}


This article advances a structural claim: governability in hybrid human–AI systems is not secured by performance, optimization, or regulatory compliance. It depends upon the sustained integrity of three irreducible structural conditions — Field, Coherence, and Limit.

These elements were not treated as analytical perspectives or policy preferences, but as mutually constitutive conditions of possibility. Field renders agency relationally intelligible. Coherence sustains its temporal compatibility across sequences of action. Limit differentiates authority and preserves the attributability of responsibility. Together, they articulate the minimal architecture within which power can remain accountable rather than merely effective.

Remove Field, and agency dissolves into operations without shared orientation.

Remove Coherence, and systemic dynamics fragment despite continued functionality.

Remove Limit, and authority persists without answerability.

The triadic structure therefore identifies the minimal architecture under which governability remains possible in systems characterized by distributed and hybridized agency. Dyadic or modular approaches may isolate necessary dimensions, yet when treated independently they permit forms of structural drift that preserve performance while eroding accountability.

The argument carries a deeper implication. Governability is not reducible to control, coordination, or alignment. It is a structural achievement: the sustained preservation of relational intelligibility, temporal compatibility, and differentiated authority. A system may scale efficiently, optimize successfully, and comply procedurally while progressively losing the very conditions that render power political. Performance without structural integrity produces fragility; compliance without relational legibility produces simulation.

Institutional responsibility follows from this architecture. Governance must assume stewardship not merely of behavior, but of the structural conditions that render governing possible. Relational configurations must remain intelligible. Systemic dynamics must remain temporally compatible. Boundaries of authority must remain differentiated and revisable. These are not enhancements to governance; they are its structural expression.

The triadic model therefore offers more than diagnostic clarity. It reorients governance toward its own conditions of possibility. It asks not only how systems perform, but whether the structural grammar that makes performance accountable remains intact.

\subsection{Absolute Closing Paragraph}
Governability is not the effect of power successfully exercised. It is the structural condition under which power can appear as political at all — locatable within a shared relational horizon, sustainable across time without self-negation, and differentiated such that responsibility remains attributable. The triadic architecture does not describe a preferred arrangement of governance. It identifies the minimal structural grammar of political order under conditions of distributed and hybridized agency.

Where these conditions endure, power remains accountable. Where they erode, governance persists only as operation or appearance. The enduring question is therefore not how to perfect control, but how to preserve the structural conditions under which control can remain answerable.

The triadic architecture is minimal because every phenomenon of governance ultimately manifests as relational configuration, temporal compatibility, or boundary differentiation.


\bibliographystyle{plainnat}
\bibliography{references}

\section*{Statements and Declarations}

\noindent\textbf{Funding} No funding was received to assist with the preparation of this manuscript.

\noindent\textbf{Competing Interests} The author declares no competing interests.

\noindent\textbf{Author Contributions} Single-author manuscript: conceptualization, theoretical development, and writing were performed by the author.

\noindent\textbf{Data Availability} No datasets were generated or analyzed during the current study.

\noindent\textbf{Code Availability} Not applicable.

\noindent\textbf{Ethics Approval} Not applicable.

\noindent\textbf{Consent to Participate} Not applicable.

\noindent\textbf{Consent for Publication} Not applicable.

\end{document}
